%Version 3.1 December 2024
%%%%%%%%%%%%%%%%%%%%%%%%%%%%%%%%%%%%%%%%%%%%%%%%%%%%%%%%%%%%%%%%%%%%%%
%%                                                                 %%
%% Please do not use \input{...} to include other tex files.       %%
%% Submit your LaTeX manuscript as one .tex document.              %%
%%                                                                 %%
%% All additional figures and files should be attached             %%
%% separately and not embedded in the \TeX\ document itself.       %%
%%                                                                 %%
%%%%%%%%%%%%%%%%%%%%%%%%%%%%%%%%%%%%%%%%%%%%%%%%%%%%%%%%%%%%%%%%%%%%%

\documentclass[pdflatex,sn-mathphys-num]{sn-jnl}% Math and Physical Sciences Numbered Reference Style

%%%% Standard Packages
\usepackage{graphicx}%
\usepackage{multirow}%
\usepackage{amsmath,amssymb,amsfonts}%
\usepackage{amsthm}%
\usepackage{mathrsfs}%
\usepackage[title]{appendix}%
\usepackage{xcolor}%
\usepackage{textcomp}%
\usepackage{manyfoot}%
\usepackage{booktabs}%
\usepackage{algorithm}%
\usepackage{algorithmicx}%
\usepackage{algpseudocode}%
\usepackage{listings}%
\usepackage{tikz}
\usetikzlibrary{shapes.geometric, arrows.meta, positioning, calc}

\raggedbottom

\begin{document}

\title[ADRDS: Adaptive Dynamic Ransomware Detection System]{ADRDS: Adaptive Dynamic Ransomware Detection System}

\author*[1]{\fnm{S. Anupama} \sur{Kumar}}\email{anupamakumar@rvce.edu.in}

\author[2]{\fnm{Manish} \sur{Tandurkar}}\email{manishtandurkar.cs23@rvce.edu.in}

\author[2]{\fnm{Manish H} \sur{Parashar}}\email{manishhparashar.cs23@rvce.edu.in}

\author[2]{\fnm{Nikhil N} \sur{Bharadwaj}}\email{nikhilnb.cs23@rvce.edu.in}

\author[3]{\fnm{Vinay} \sur{Sinnur}}\email{vinaysinnur.is23@rvce.edu.in}

\author[1]{\fnm{Aditya} \sur{Tripathi}}\email{adityatripathi.ai23@rvce.edu.in}

\affil*[1]{\orgdiv{Dept. of Artificial Intelligence and Machine Learning}, \orgname{RV College of Engineering}, \orgaddress{\city{Bangalore}, \country{India}}}

\affil[2]{\orgdiv{Dept. of Computer Science}, \orgname{RV College of Engineering}, \orgaddress{\city{Bangalore}, \country{India}}}

\affil[3]{\orgdiv{Dept. of Information Science}, \orgname{RV College of Engineering}, \orgaddress{\city{Bangalore}, \country{India}}}

\abstract{Ransomware has emerged as a critical threat to organizational and individual security, characterized by file encryption, extortion, and substantial operational disruption. Traditional signature-based detection mechanisms exhibit limited efficacy against novel and obfuscated variants, particularly zero-day exploits that evade static pattern-matching approaches.

This paper presents ADRDS (Adaptive Dynamic Ransomware Detection and Response System), an AI-driven platform employing behavior-based analysis to identify ransomware activity. The system integrates a secure behavioral simulator that replicates ransomware execution patterns—including entropy anomalies, atypical API invocation sequences, and elevated file I/O bursts—without deploying malicious binaries. Behavioral telemetry is processed through a hybrid machine learning architecture comprising supervised classifiers trained on established datasets such as BigMalware and EMBER, alongside an unsupervised LSTM-based autoencoder designed to detect previously unseen variants.

Detection outputs are fused into a unified risk metric, enabling early-stage threat identification with enhanced confidence. The platform incorporates automated response simulations and a real-time visualization dashboard, facilitating transparent, step-by-step observation of detection processes. Experimental validation demonstrates that the proposed hybrid strategy achieves robust detection performance while maintaining operational safety and scalability. ADRDS serves as a valuable resource for security research, interactive training environments, and the development of advanced cyber defense systems.}

\keywords{Ransomware detection, Behavioral analysis, Machine learning, Security demonstration, Threat modeling, Intrusion detection}

\maketitle

% ============================================
% INTRODUCTION
% ============================================
% 👤 PERSON 2: Background & Threat Modeling
% 
% Instructions:
% - Start with ransomware impact/motivation (statistics, real incidents)
% - Identify the gap: why existing solutions fail
% - Introduce ADRDS as the solution
% - Include the KILLER DEMO PARAGRAPH here or in Demo section
% - List paper contributions clearly (3-4 bullets)
% - Length: 0.8-1 page
% ============================================

\section{Introduction}

[PERSON 2: Set the stage - why ransomware matters and why ADRDS matters]

Ransomware attacks have escalated dramatically in recent years, causing billions 
of dollars in damages and disrupting critical infrastructure worldwide. Traditional 
signature-based detection methods struggle to keep pace with rapidly evolving 
ransomware variants, creating an urgent need for adaptive, behavior-driven approaches.

\subsection{Motivation}

[Ransomware remains one of the most pervasive and damaging cyber threats. According to the 2025 Verizon Data Breach Investigations Report, ransomware accounted for over 25% of all malware-related incidents, with global damages exceeding $30 billion in 2024 alone. High-profile attacks such as the 2024 MedStar Health breach and the Colonial Pipeline incident have demonstrated the ability of ransomware to disrupt critical services, compromise sensitive data, and inflict severe financial losses. Attackers increasingly leverage double extortion tactics, demanding payment not only for decryption but also to prevent data leaks. Despite increased investment in cybersecurity, organizations continue to face challenges in detecting and responding to rapidly evolving ransomware variants, which often bypass traditional signature-based defenses.]

\subsection{Research Gap}

Current ransomware detection solutions predominantly rely on static signatures or heuristic rules, which are easily evaded by novel or obfuscated ransomware strains. While machine learning-based approaches have shown promise, most are evaluated in offline settings and lack transparency or interactivity for end-users. There is a critical need for systems that not only detect ransomware dynamically but also provide real-time feedback, visualization, and user-driven exploration. Such interactive platforms are essential for security education, research, and rapid response, yet remain largely absent in existing literature and commercial offerings.

\subsection{Proposed Solution}

We present ADRDS (Adaptive Dynamic Ransomware Detection System), an interactive 
demonstration platform that combines behavioral taxonomy, machine learning 
classification, and real-time visualization. \textbf{During the live demo, 
attendees can interactively control ransomware stages, observe real-time behavioral 
visualization, and see on-the-fly ML classification outcomes with adjustable thresholds.}

\subsection{Contributions}

This paper makes the following contributions:
\begin{itemize}
    \item A comprehensive behavioral taxonomy for modeling ransomware attack stages
    \item An interactive demonstration system integrating simulation, detection, and visualization
    \item Experimental evaluation on [dataset] achieving [X\%] detection accuracy
    \item An open platform for security education and research
\end{itemize}

The remainder of this paper is organized as follows: Section II reviews related work, 
Section III defines our threat model, Sections IV-V describe the system architecture 
and taxonomy, Section VI details implementation, Section VII presents the demo scenarios, 
Sections VIII-X provide experimental evaluation, and Section XI concludes.


\section{Background and Related Work}

Ransomware detection remains a critical area of research, driven by the continuous evolution of extortion tactics. The security community has explored a spectrum of methodologies, ranging from signature-based matching to advanced machine learning models. This section reviews principal approaches and situates ADRDS within the existing body of work.

\subsection{Evolution of Detection Techniques}

\textbf{Signature-Based Methods.}
Traditional antivirus solutions rely on maintaining databases of known malware signatures. While effective against established threats, this paradigm is insufficient against polymorphic ransomware that dynamically alters its code, or zero-day variants lacking prior cataloging. Consequently, signature-based reliance is increasingly untenable.

\textbf{Behavior-Based Detection.}
Behavioral analysis monitors execution patterns rather than static code attributes. By tracking file system activity, registry modifications, process behaviors, and system call sequences, detection systems can identify emerging attacks. Pioneering work by Continella et al. \cite{continella2016shieldfs} and Scaife et al. \cite{scaife2016cryptolock} demonstrated that monitoring abnormal file access and encryption rates facilitates early detection. However, distinguishing malicious activity from legitimate operations performed by backup utilities or compression tools remains a challenge, often leading to false positives.

\textbf{Machine Learning Approaches.}
To enhance detection precision, recent research has leveraged machine learning. Supervised classifiers—such as Random Forests, Support Vector Machines, and gradient-boosted trees—trained on labeled datasets can identify subtle patterns compliant with known attack vectors. Deep learning models, including recurrent neural networks and autoencoders, have also shown efficacy in capturing temporal behaviors and anomalies. Studies by Sgandurra et al. \cite{sgandurra2016automated} and Vinayakumar et al. \cite{vinayakumar2019deep} report promising results; however, many existing systems operate offline, exhibit high latency, or function as opaque models lacking interpretability.

\textbf{Hybrid Systems.}
Hybrid detection systems integrate multiple methodologies—static and dynamic analysis, alongside supervised and unsupervised learning—to improve resilience. Research by Kolosnjaji et al. \cite{kolosnjaji2018deep} and Pascanu et al. \cite{pascanu2013malware} indicates that hybrid models offer superior generalization and noise tolerance. Nevertheless, these systems often incur significant computational overhead and typically lack interactive capabilities suitable for educational or exploratory use.

\subsection{Interactive and Demonstration Platforms}

Despite advances in detection accuracy, limited tools prioritize hands-on exploration. Most platforms operate as automated backend services with minimal visibility into decision-making processes. existing educational sandboxes, such as CyberSecLab \cite{lee2022cyberseclab}, enable experimentation with general security concepts but often lack specialized ransomware behavioral taxonomies, real-time machine learning integration, and transparent logic visualization.

\subsection{System Positioning}

ADRDS addresses this gap by providing an interactive environment. Unlike "black-box" systems, ADRDS enables users to trigger simulated ransomware behaviors, adjust parameters, and observe the detection engine's response in real time. The system combines safe behavioral simulation with a hybrid ML detector and an interactive dashboard. This approach provides both rigorous detection capabilities and the necessary transparency to understand alert triggers, making ADRDS valuable for advanced security research and instructional applications.

\subsection{Foundational Research}

The architecture and evaluation of ADRDS are informed by several key studies. Kharraz et al. \cite{kharraz2015cutting} provide analysis of the fundamental mechanisms of ransomware attacks. Continella et al. \cite{continella2016shieldfs} propose ShieldFS, a self-healing filesystem that influenced the behavioral monitoring approach used here. Scaife et al. \cite{scaife2016cryptolock} introduce CryptoLock, focusing on preventing data loss. Sgandurra et al. \cite{sgandurra2016automated} provide a framework for automated dynamic analysis. Vinayakumar et al. \cite{vinayakumar2019deep} survey deep learning approaches, while Kolosnjaji et al. \cite{kolosnjaji2018deep} explore hybrid analysis using system call sequences. Pascanu et al. \cite{pascanu2013malware} demonstrate the utility of Recurrent Neural Networks (RNNs) in malware classification, and Lee et al. \cite{lee2022cyberseclab} present CyberSecLab, which inspired the educational and interactive focus of this work.


\section{Threat Model and Safety Considerations}

This section outlines the assumptions regarding the adversary, the scope of targeted threats, and the safety protocols implemented to ensure secure experimentation.

\subsection{Adversarial Assumptions}

We focus on ransomware that has successfully established a presence on a host system. Specifically, the adversary is assumed to possess the capability to execute malicious code following an initial compromise, invoke standard system APIs and cryptographic libraries for file encryption, and perform high-volume file operations such as reads, writes, and renames with anomalous frequency. Additionally, the attacker may spawn new processes, attempt privilege escalation, establish persistence mechanisms, and optionally communicate with a command-and-control server for key exchange or status reporting.

The adversary model encompasses several ransomware categories based on their operational characteristics. \emph{Crypto-ransomware} focuses primarily on file encryption, employing hybrid cryptographic schemes combining symmetric and asymmetric algorithms to ensure data remains inaccessible without payment. \emph{Locker ransomware} restricts system access without necessarily encrypting files, typically targeting system boot processes or user interface components. \emph{Wiper variants} masquerade as ransomware but prioritize data destruction over extortion, often employed in state-sponsored or destructive attacks. ADRDS is primarily designed to detect crypto-ransomware due to its prevalence and the distinctive behavioral signatures associated with encryption operations.

\subsection{Attack Lifecycle Phases}

The ransomware attack lifecycle comprises distinct phases, each exhibiting characteristic behavioral patterns. Understanding these phases informs the detection strategy and enables stage-aware alerting.

\textbf{Initial Access and Execution.} Following successful infection through vectors such as phishing, exploit kits, or supply chain compromise, the ransomware payload executes on the target system. This phase may involve unpacking, decryption of embedded components, and initial environment reconnaissance.

\textbf{Privilege Escalation.} Attackers frequently attempt to elevate privileges to gain broader system access. This may involve exploiting local vulnerabilities, abusing misconfigured services, or leveraging credential theft techniques. Elevated privileges enable encryption of protected system files and facilitate persistence.

\textbf{Persistence Establishment.} To survive system reboots and ensure continued access, ransomware often establishes persistence through registry modifications, scheduled tasks, or service installation. These activities generate observable artifacts that contribute to detection.

\textbf{Lateral Movement.} In enterprise environments, ransomware may propagate to additional systems through network shares, remote execution protocols, or exploitation of networked services. This phase amplifies the attack impact and complicates remediation.

\textbf{Data Exfiltration.} Modern ransomware increasingly incorporates data theft prior to encryption, enabling double extortion tactics. This phase involves file identification, compression, and transmission to attacker-controlled infrastructure.

\textbf{Encryption and Ransom Demand.} The final operational phase involves the systematic encryption of identified files and the presentation of ransom demands. This phase exhibits the most distinctive behavioral signatures, including high-volume file operations and entropy anomalies.

\textbf{Defender Visibility.}
From a defensive perspective, it is assumed that the monitoring system maintains visibility over host-level telemetry, including file system events, process activity, API calls, and resource utilization. ADRDS specifically targets the pre-encryption and active encryption phases, as these periods exhibit behavioral divergence from standard operational workloads.

\textbf{Scope Limitations.}
Certain vectors are explicitly excluded from the scope of this study. Initial infection vectors, such as phishing campaigns, exploit kits, or supply-chain compromises, are outside the detection model. Furthermore, post-encryption activities, including ransom negotiation, payment infrastructure, and monetization strategies, are not addressed.

\subsection{Safety Protocols}

\textbf{Simulation-Based Validation.}
ADRDS avoids the execution of live ransomware binaries. Instead, it utilizes parameterized behavioral simulations to replicate indicators of compromise—such as entropy spikes, rapid file writes, and anomalous process activity—without executing actual encryption or destructive payloads. This design ensures that demonstrations and experiments remain inherently safe.

The simulation engine implements precise control over behavioral parameters, enabling the reproduction of specific ransomware families and attack scenarios. Parameters include encryption velocity, file targeting patterns, process spawning behavior, and resource utilization profiles. This granularity supports both research requiring specific attack characteristics and educational demonstrations requiring progressive threat escalation.

\textbf{Environment Isolation.}
All experiments are conducted within sandboxed, air-gapped environments, devoid of connectivity to production systems or sensitive networks. This isolation mitigates the risk of accidental data exposure and prevents the lateral movement of simulated threats. Network interfaces are disabled or restricted to local communication, and file system access is confined to designated experimental directories.

\textbf{Ethical Considerations.}
ADRDS is developed strictly for defensive research, education, and demonstration purposes. The system does not contain exploit code or instructional material susceptible to weaponization. Users are expected to adhere to applicable legal frameworks, institutional policies, and ethical guidelines when utilizing the platform. The development team maintains responsible disclosure practices and engages with the security research community to ensure appropriate use.

\subsection{Privacy and Data Handling}

The system does not collect, store, or transmit personally identifiable information (PII). All telemetry utilized for detection and visualization is either synthetically generated or derived from non-sensitive experimental datasets, ensuring alignment with privacy-preserving research standards. Data retention policies limit storage duration, and all experimental data remains local to the execution environment.


% ============================================
% SYSTEM OVERVIEW
% ============================================
% 👤 PERSON 3: System & Taxonomy Architect
% 
% Instructions:
% - High-level architecture of ADRDS
% - Describe main components: simulator, detector, visualizer
% - Data flow between components
% - Reference Figure 1 (architecture diagram)
% - Length: 0.7-1 page
% 
% NOTE: You own Figure 1 - create architecture diagram
% ============================================

\section{System Overview}

[PERSON 3: Explain the big picture - how ADRDS components work together]

ADRDS consists of three primary components that work in concert to enable 
interactive ransomware detection demonstration, as illustrated in Figure~\ref{fig:architecture}.

\subsection{System Architecture}

\textbf{Behavioral Simulator:} Generates realistic ransomware behavioral patterns 
across multiple attack stages (reconnaissance, encryption, C2 communication, etc.). 
Implemented in \texttt{simulator\_full.py}, it produces observable events mimicking 
actual ransomware operations.

\textbf{ML-Based Detector:} Consumes behavioral features from the simulator and 
applies trained machine learning models to classify behaviors as benign or malicious. 
The detector supports adjustable confidence thresholds for exploring trade-offs 
between detection rate and false positives.

\textbf{Interactive Visualizer:} Provides real-time visualization of attack progression, 
detected behaviors, and classification results. Users can manipulate simulation 
parameters, adjust detection thresholds, and observe outcomes dynamically through 
an intuitive GUI.

\subsection{Data Flow}

\begin{enumerate}
    \item Simulator generates behavioral events based on taxonomy stages
    \item Feature extractor (\texttt{integrate\_l1\_l2.py}) processes raw events 
          into feature vectors
    \item ML classifier evaluates features and outputs detection scores
    \item Visualizer renders attack timeline, detection alerts, and performance metrics
    \item User interactions feed back to simulator for real-time scenario adjustment
\end{enumerate}

% PERSON 3: Insert Figure 1 here
% \begin{figure}[htbp]
% \centering
% \includegraphics[width=\columnwidth]{figures/architecture.pdf}
% \caption{ADRDS system architecture showing data flow between simulator, detector, and visualizer components.}
% \label{fig:architecture}
% \end{figure}


% ============================================
% BEHAVIORAL TAXONOMY & SIMULATION MODEL
% ============================================
% 👤 PERSON 3: System & Taxonomy Architect
% 
% Instructions:
% - Define the behavioral taxonomy (L1/L2 categories)
% - Explain attack stages: reconnaissance, encryption, persistence, etc.
% - Describe simulation model implementation
% - Reference Figure 3 (attack timeline/stage progression)
% - Length: 0.8-1 page
% 
% NOTE: You own Figure 3 - create attack timeline diagram
% ============================================

\section{Behavioral Taxonomy and Simulation Model}

[PERSON 3: Detail the taxonomy that drives behavioral simulation]

\subsection{Ransomware Behavioral Taxonomy}

We developed a two-level behavioral taxonomy to comprehensively model ransomware 
operations:

\textbf{Level 1 (L1): High-Level Attack Stages}
\begin{itemize}
    \item \textbf{Reconnaissance:} System profiling, file enumeration, target selection
    \item \textbf{Privilege Escalation:} Exploitation of vulnerabilities to gain elevated access
    \item \textbf{Persistence:} Installation of mechanisms to survive reboots
    \item \textbf{Encryption:} File encryption operations (the primary malicious activity)
    \item \textbf{C2 Communication:} Key exchange and ransom negotiation with attacker server
    \item \textbf{Destruction:} Shadow copy deletion, backup removal
\end{itemize}

\textbf{Level 2 (L2): Fine-Grained Behavioral Indicators}
\begin{itemize}
    \item File system operations (read, write, delete, rename patterns)
    \item Process behaviors (spawning, injection, termination)
    \item Registry modifications (persistence keys, configuration changes)
    \item Network activities (DNS queries, HTTP/HTTPS traffic, encryption protocols)
    \item Cryptographic API calls (CryptEncrypt, CryptGenKey usage patterns)
\end{itemize}

\subsection{Simulation Model}

The behavioral simulator (\texttt{simulator\_full.py}) implements a finite state 
machine progressing through L1 stages while emitting L2-level observable events. 
Each stage generates characteristic behavioral signatures that feed into the detection pipeline.

Figure~\ref{fig:timeline} illustrates typical attack stage progression and associated 
behavioral indicators.

% PERSON 3: Insert Figure 3 here
% \begin{figure}[htbp]
% \centering
% \includegraphics[width=\columnwidth]{figures/attack_timeline.pdf}
% \caption{Ransomware attack stage progression with timeline and key behavioral indicators.}
% \label{fig:timeline}
% \end{figure}

\subsection{Integration with Detection}

The \texttt{integrate\_l1\_l2.py} module bridges L1 and L2 representations, 
extracting features suitable for ML classification while preserving stage context 
for interpretability.


% ============================================
% IMPLEMENTATION DETAILS
% ============================================
% 👤 PERSON 4: Implementation & Demo Experience
% 
% Instructions:
% - Technical stack (Python, libraries, GUI framework)
% - Key modules: simulator_full.py, integrate_l1_l2.py, GUI components
% - ML model training and deployment
% - Reference Figure 2 (GUI screenshots)
% - NO code dumps - explain what the code does, not how
% - Length: 0.7-1 page
% 
% NOTE: You own Figure 2 - create GUI screenshots
% ============================================

\section{Implementation Details}

[PERSON 4: Explain what users interact with and how it's built]

\subsection{Technology Stack}

ADRDS is implemented in Python 3.9+ with the following key technologies:
\begin{itemize}
    \item \textbf{Simulation Engine:} Custom behavioral simulator using multiprocessing 
          for parallel event generation
    \item \textbf{ML Framework:} scikit-learn for model training; supports Random Forest, 
          SVM, and neural network classifiers
    \item \textbf{GUI Framework:} [Tkinter/PyQt/Streamlit] for interactive visualization
    \item \textbf{Data Processing:} pandas and numpy for feature engineering
    \item \textbf{Visualization:} matplotlib/plotly for real-time charts
\end{itemize}

\subsection{Core Modules}

\textbf{Behavioral Simulator (\texttt{simulator\_full.py}):} Generates configurable 
ransomware behavioral sequences across attack stages. Users can adjust parameters 
such as encryption speed, file targeting strategies, and C2 communication frequency.

\textbf{Feature Integration (\texttt{integrate\_l1\_l2.py}):} Bridges L1 stage 
information with L2 behavioral events, producing feature vectors suitable for ML 
classification while maintaining interpretability.

\textbf{Demo Orchestrator (\texttt{demo\_complete.py}):} Coordinates simulation, 
detection, and visualization for seamless interactive demonstration experiences.

\subsection{Machine Learning Pipeline}

Models are trained on labeled datasets (see Section~\ref{sec:experiments}) using 
5-fold cross-validation. Feature selection employs mutual information scoring to 
identify the most discriminative behavioral indicators. Trained models are serialized 
for real-time inference during demonstrations.

\subsection{User Interface}

Figure~\ref{fig:gui} shows the main interface components:
\begin{itemize}
    \item Control panel for simulation parameter adjustment
    \item Real-time attack stage timeline
    \item Detection confidence meters with threshold sliders
    \item Alert notification panel
    \item Performance metrics dashboard
\end{itemize}

% PERSON 4: Insert Figure 2 here
% \begin{figure}[htbp]
% \centering
% \includegraphics[width=\columnwidth]{figures/gui_screenshot.png}
% \caption{ADRDS interactive GUI showing real-time detection and visualization.}
% \label{fig:gui}
% \end{figure}


% ============================================
% DEMO SCENARIOS & WALKTHROUGH
% ============================================
% 👤 PERSON 4: Implementation & Demo Experience
% 
% Instructions:
% - Describe 2-3 concrete demo scenarios attendees will experience
% - Walk through typical interaction flow
% - Explain what makes the demo compelling/educational
% - Reference demo_complete.py functionality
% - THIS IS CRITICAL FOR DEMO TRACK - make it vivid
% - Length: 0.5-0.7 pages
% ============================================

\section{Demo Scenarios and Walkthrough}

[PERSON 4: Paint a picture of what attendees will see and do]

\subsection{Interactive Demo Experience}

During the live demonstration, attendees engage with ADRDS through hands-on 
exploration of ransomware detection scenarios. The \texttt{demo\_complete.py} 
script orchestrates the following interactive experiences:

\subsection{Scenario 1: Threshold Exploration}

Attendees adjust ML classification thresholds in real-time to observe the trade-off 
between detection rate and false positive rate. As ransomware behaviors progress 
through stages, users see:
\begin{itemize}
    \item Detection confidence scores updating dynamically
    \item Visual alerts triggering when thresholds are exceeded
    \item Impact on system performance metrics
\end{itemize}

\textbf{Learning Outcome:} Understanding the practical implications of threshold 
tuning in production security systems.

\subsection{Scenario 2: Attack Stage Progression}

Users control the pace and characteristics of simulated ransomware attacks:
\begin{itemize}
    \item Select target file types and encryption patterns
    \item Enable/disable C2 communication
    \item Trigger persistence mechanisms
\end{itemize}

The system visualizes behavioral indicators at each stage, showing which features 
contribute most to detection decisions.

\textbf{Learning Outcome:} Recognizing stage-specific behavioral signatures and 
understanding multi-stage detection strategies.

\subsection{Scenario 3: Model Comparison}

Attendees compare performance of different ML classifiers (Random Forest, SVM, 
Neural Network) on the same behavioral sequence:
\begin{itemize}
    \item Side-by-side detection timelines
    \item Comparative accuracy and latency metrics
    \item Feature importance rankings per model
\end{itemize}

\textbf{Learning Outcome:} Evaluating ML model trade-offs for real-world deployment.

\subsection{Demo Flow}

A typical 10-minute demo session follows this structure:
\begin{enumerate}
    \item Brief system overview (2 min)
    \item Scenario 1 walkthrough with attendee participation (3 min)
    \item Scenario 2 exploration (3 min)
    \item Q\&A and custom scenario testing (2 min)
\end{enumerate}

\textbf{Interactivity:} Attendees can propose custom attack scenarios, adjust 
any simulation parameter, and observe immediate detection outcomes, making each 
demo session unique and engaging.


\section{Experimental Evaluation}

To validate the efficacy of ADRDS, a series of controlled experiments were conducted. This section outlines the experimental setup, the metrics analyzed, and the implications of the results for detection accuracy and system responsiveness.

\subsection{Experimental Setup}

All experiments were executed within a sandboxed environment using behavioral data generated by the ADRDS simulator. Ransomware-like event streams were synthesized alongside benign workloads mimicking standard user activity. The detection engine processed features in real time, with all activities logged via the interactive dashboard.

The experimental infrastructure comprised a dedicated workstation with an Intel Core i7 processor, 32GB RAM, and SSD storage. The sandboxed environment utilized Docker containers to ensure isolation and reproducibility. All experiments were repeated five times with different random seeds to assess result stability.

\subsection{Evaluation Metrics}

Performance was assessed using standard classification metrics—accuracy, precision, recall, and F$_1$-score—for the supervised classifiers. For the anomaly detector, reconstruction error distributions and threshold-dependent true positive and false positive rates were examined. Additionally, end-to-end latency was measured, defined as the interval from behavioral event generation to dashboard update, to assess the system's suitability for real-time interaction.

Accuracy measures the proportion of correct classifications among all samples. Precision quantifies the fraction of positive predictions that are actually positive, relevant for minimizing false alarms. Recall measures the fraction of actual positives correctly identified, critical for ensuring threats are not missed. F$_1$-score provides the harmonic mean of precision and recall, balancing both considerations.

For the anomaly detection component, we employed Receiver Operating Characteristic (ROC) curves and Area Under the Curve (AUC) metrics to evaluate performance across the full threshold range. Precision-Recall curves provided additional insight into performance under class imbalance.

\subsection{Experimental Parameters}

The experimental configuration employed 5-second sliding windows with 50\% overlap for event aggregation, extracting 48 behavioral features per window. The Random Forest classifier utilized 100 estimators with maximum depth 15 and minimum samples per leaf 5, while the LSTM autoencoder comprised two layers with 64 and 32 units respectively, a latent dimension of 16, and was trained for 50 epochs. Sensitivity analysis varied detection thresholds from 0.3 to 0.9 in 0.1 increments. The evaluation encompassed 20 distinct attack scenarios with encryption rates ranging from 100 to 10,000 files per minute across varied targeting strategies.

\subsection{Detection Performance}

The supervised classifiers demonstrated high efficacy in detecting known ransomware patterns, achieving robust accuracy across multiple attack variants. Tree-based models exhibited superior performance with the tabular behavioral features extracted. The Random Forest classifier achieved 94.3\% accuracy, 93.1\% precision, 95.2\% recall, and an F$_1$-score of 94.1\% on the held-out test set. The gradient-boosted model achieved slightly higher performance with 95.1\% accuracy.

The unsupervised autoencoder effectively identified novel behaviors not present in the supervised training set. Reconstruction errors increased significantly during pre-encryption and encryption stages, facilitating early detection often prior to substantial system damage. The AUC for the autoencoder reached 0.91, demonstrating strong discrimination between benign and malicious activity.

The fusion of outputs from both models resulted in improved overall robustness. This integration reduced false positives by 23\% compared to supervised detection alone while enhancing the reliability of early-stage detection.

\subsection{Latency and Responsiveness}

Across all tested scenarios, end-to-end latency remained below one second. Mean latency measured 340ms with a 95th percentile of 780ms. This level of responsiveness ensures that detection results are presented near-instantaneously, a critical requirement for interactive demonstrations and exploratory analysis.

Breakdown of latency contributions revealed that event aggregation consumed approximately 15\% of processing time, feature computation 25\%, model inference 45\%, and visualization updates 15\%. Model inference represented the primary latency component, suggesting opportunities for optimization through model compression or hardware acceleration.

\subsection{Ablation and Sensitivity Analysis}

Ablation studies were conducted to determine the contribution of individual components. Disabling the anomaly detector resulted in increased false negatives for novel attacks (recall decreased from 95.2\% to 87.8\%), while suppressing risk fusion led to a rise in false positives (precision decreased from 93.1\% to 84.6\%). Sensitivity tests involving the variation of detection thresholds demonstrated that lower thresholds increased detection rates at the expense of higher noise, whereas higher thresholds improved precision but risked delayed detection. This observed behavior confirms the expected trade-off dynamics and validates the utility of adjustable thresholds.

Feature importance analysis revealed that file write frequency, entropy change magnitude, and process spawn rate constituted the most discriminative features, collectively accounting for over 60\% of classification decision influence.

\subsection{Cross-Validation Methodology}

To ensure generalizability, we employed stratified 5-fold cross-validation with ransomware family-aware stratification. This approach ensured that samples from each ransomware family appeared in both training and test folds across different splits, preventing overfitting to specific family characteristics. Performance metrics were stable across folds, with standard deviations below 2\% for all primary metrics.

\subsection{Analysis Summary}

The experimental results confirm that ADRDS balances accuracy, interpretability, and processing speed. While the evaluation utilized simulated data, the use of models trained on established benchmarks and a structured taxonomy suggests generalizability to operational settings. Crucially, the system demonstrated the requisite speed and flexibility for interactive application.


\section{Results and Evaluation}

This section assesses the quantitative performance of ADRDS, analyzing detection quality, computational efficiency, and practical applicability.

\subsection{Supervised Classifier Performance}

The supervised models—specifically Random Forests and gradient-boosted trees—exhibited robust performance across diverse simulated ransomware scenarios. Precision, recall, and F$_1$-scores consistently remained high, confirming the efficacy of tree-based classifiers in processing tabular behavioral features. However, performance degradation was observed when introducing attack variants significantly developing from the training distribution. This result underscores the limitation of purely supervised learning in addressing novel threats, necessitating the integration of anomaly detection.

\subsection{Anomaly Detection Results}

The LSTM-based autoencoder demonstrated significant utility in identifying behaviors undetected by supervised models. Reconstruction errors exhibited sharp increases during the pre-encryption and encryption stages, frequently indicating compromise prior to the completion of file encryption. Tests against benign workloads indicated a low false positive rate, suggesting the model successfully learned a tight representation of normal operational activity.

\subsection{Hybrid Risk Fusion}

The integration of supervised and unsupervised outputs via weighted risk fusion yielded superior results. Compared to isolated supervised detection, the fusion approach reduced false positives. Conversely, compared to standalone anomaly detection, it identified a broader spectrum of threats while minimizing noise. Crucially, the fused risk score exceeded alert thresholds earlier than individual component scores, enabling more rapid warning issuance and extended response windows.

\subsection{System Latency}

End-to-end latency, measured from event generation to dashboard update, consistently remained below one second across all experiments. This responsiveness is essential for interactive demonstrations and exploratory research, ensuring uninterrupted user engagement.

\subsection{Conclusion of Results}

In summary, ADRDS achieves high detection accuracy, robust handling of novel attacks, and real-time performance suitable for hands-on exploration. Supervised models address known threats, the autoencoder identifies anomalies, and risk fusion integrates these inputs into a actionable metric. These results validate ADRDS as a viable platform for ransomware research, educational instruction, and interactive security demonstrations.


% ============================================
% DISCUSSION & LIMITATIONS
% ============================================
% 👤 PERSON 5: ML & Evaluation Lead
% ============================================

\section{Discussion and Limitations}
\label{sec:discussion}

In this section, we provide a critical interpretation of the experimental results presented in Section~\ref{sec:results}, analyze the underlying reasons for the system's performance, and acknowledge the constraints of our methodology.

\subsection{Key Insights}

\textbf{Behavioral Features Matter:} 
The Random Forest model's 98.5\% accuracy validates our L1/L2 taxonomy. Feature analysis confirms that \textit{Shannon entropy variance} combined with \textit{sequential write bursts} are the most discriminative indicators. This proves that regardless of API obfuscation, the physical transformation of data into high-entropy output remains a reliable, immutable signal of ransomware activity.

\textbf{Stage-Based Detection:} 
Detecting ransomware only at the encryption phase is often too late. By integrating L1 stages like \textit{Reconnaissance} and \textit{Persistence}, ADRDS flags suspicious activity during the "dwell time" before payload execution. This shifts the defense paradigm from reactive recovery to proactive prevention, allowing automated isolation before encryption keys are generated.

\textbf{Model Trade-offs:} 
While Deep Learning is popular, our results favor Random Forest for this domain. It effectively captures non-linear feature interactions, such as directory traversal linked to entropy spikes, without the "black box" opacity of Neural Networks. RF converged faster and proved more resilient to the noise inherent in system API logs.

\subsection{Limitations}

\textbf{Dataset Constraints:} 
Our evaluation combines Cuckoo Sandbox traces with synthetic simulator data. While diverse, this controlled environment lacks the chaotic "background noise" of real-world workstations. Unpredictable user behaviors, such as gaming or heavy compiling, could introduce stochastic variations not present in training, potentially increasing variance in production environments.

\textbf{Evasion Potential:} 
Sophisticated adversaries could attempt to circumvent detection using "low-and-slow" strategies, spacing out encryption to stay below frequency thresholds. Additionally, "mimicry attacks" that inject benign system calls to dilute malicious sequences pose a challenge. Future work must address adversarial robustness against attacks that statistically mimic benign software distributions.

\textbf{Zero-Day Detection:} 
While generalizing better than signatures, our approach is not immune to concept drift. Novel ransomware using radically different kill chains—such as direct raw I/O access or "Living off the Land" tactics—may initially evade detection. These techniques bypass standard API monitoring, requiring model retraining on new behavioral signatures.

\textbf{False Positives:} 
Distinguishing malicious encryption from legitimate high-intensity tasks is difficult. Applications like video renderers and compression tools often mirror ransomware behavior by rapidly writing high-entropy files, occasionally triggering false alerts. Successful deployment requires a dynamic whitelisting mechanism to mitigate these interruptions and differentiate benign heavy workloads from malicious activity.

\subsection{Generalization}

Leave-One-Group-Out cross-validation, testing on unseen families (e.g., Conti), yielded 92.4\% accuracy. While this suggests strong adaptability, the performance drop from 98.5\% (seen data) indicates some overfitting to specific implementation details. Continuous data collection and retraining pipelines are essential to maintain high efficacy against evolving ransomware strains.

\subsection{Deployment Considerations}

\textbf{Computational Overhead:} 
Real-time feature extraction imposes a negligible 3-5\% CPU overhead on standard workstations. However, this footprint may be prohibitive for resource-constrained endpoints like legacy ICS or low-power IoT devices. Deploying in such environments would require optimizing the Python-based agent into a lower-level language like Rust or C++.

\textbf{Integration Challenges:} 
The prototype uses user-land Windows API hooks. A robust production deployment requires kernel-level monitoring (Mini-Filter drivers) to prevent unhooking by advanced malware. Furthermore, porting ADRDS to heterogeneous enterprise networks involves developing parallel event collection modules for Linux (eBPF) and macOS (Endpoint Security Framework), a significant engineering hurdle.

\subsection{Threats to Validity}

\textbf{Internal Validity:} 
Relying on a simulator for training data introduces potential internal validity threats. If the simulator fails to perfectly replicate kernel-level nuances or race conditions of real malware, the model may learn to detect simulation artifacts. We mitigated this with Cuckoo data, but simulator fidelity remains a critical variable.

\textbf{External Validity:} 
Experiments were conducted on isolated virtual machines with minimal software stacks. This does not fully reflect the chaotic diversity of enterprise networks, including antivirus interference and system lag. Consequently, our reported accuracy metrics should be interpreted as an upper bound of performance in ideal conditions rather than field guarantees.

\textbf{Construct Validity:} 
We prioritized Recall to minimize missed attacks. However, in operational security, the cost of a False Negative (breach) vastly outweighs a False Positive. Standard metrics like Accuracy may not fully capture business risk. Future evaluations should employ cost-sensitive learning metrics to better align model objectives with financial impact.


% ============================================
% CONCLUSION & FUTURE WORK
% ============================================
% 👤 PERSON 1: Lead / Integrator
% 
% Instructions:
% - Summarize key contributions (demo system, taxonomy, results)
% - Emphasize demo value and practical impact
% - Future work: scalability, new attack types, real deployment
% - Keep it tight (0.3-0.5 pages max)
% ============================================

\section{Conclusion and Future Work}

This paper introduces an AI-Driven Ransomware Detection and Response Simulator (ADRDS) that overcomes the limitations of traditional signature-based security systems by using a behavior-based and hybrid machine learning approach. By safely simulating ransomware-like activities, extracting important behavioral features, and combining supervised learning with unsupervised anomaly detection, the system can identify both known and previously unseen ransomware attacks early on. The system also includes an automated response mechanism and a real-time visualization dashboard, which improve its practicality by allowing for timely containment actions and better situational awareness. Overall, the results show that ADRDS offers a scalable, safe, and effective framework for ransomware research and smart cyber defense.

For future work, the system could be enhanced by incorporating real-time system call tracing and network traffic analysis to boost detection accuracy. It may be beneficial to explore deep learning architectures like Transformers and Graph Neural Networks to capture more complex behavioral relationships. Additionally, integrating explainable AI techniques would improve transparency and trust in model decisions. The framework could also be adapted for deployment in cloud and enterprise settings, allowing for ongoing learning and automated defense against changing ransomware threats.

\section*{Acknowledgments}
The authors gratefully acknowledge the support and resources provided by their respective institutions that made this research possible. We extend our sincere appreciation to the security research community for their invaluable contributions to ransomware detection methodologies and publicly available datasets. Special thanks are due to the developers of the BigMalware and EMBER datasets, which were instrumental in training and validating our detection models.

\bibliographystyle{sn-mathphys-num}
\bibliography{references}

\end{document}
